\documentclass[10pt,aspectratio=169]{beamer}
\usepackage[portuguese]{babel}
\usepackage{fontspec,unicode-math}
\setmainfont{Noto Sans}
\setsansfont{Noto Sans}
\setmathfont{STIX Math}
\usepackage{amsmath}
\usepackage{xcolor}
\definecolor{InsperRed}{HTML}{a01924}
\definecolor{Ink}{HTML}{232529}
\setbeamercolor{frametitle}{fg=Ink}
\setbeamercolor{structure}{fg=InsperRed}
\setbeamercolor{normal text}{fg=Ink,bg=white}
\setbeamerfont{frametitle}{size=\large,series=\bfseries}
\setbeamerfont{title}{size=\LARGE,series=\bfseries}
\setbeamertemplate{navigation symbols}{}
\setbeamertemplate{itemize item}{\color{InsperRed}$\bullet$}
\setbeamertemplate{footline}{%
  \hbox{\begin{beamercolorbox}[wd=.84\paperwidth,ht=2.7ex,dp=1.2ex,leftskip=.7cm]{normal text}\scriptsize Inteligência Artificial para Decisão · Insper · 2027-1\end{beamercolorbox}%
  \begin{beamercolorbox}[wd=.16\paperwidth,ht=2.7ex,dp=1.2ex,rightskip=.7cm plus1fil]{normal text}\hfill\scriptsize\insertframenumber\, /\, \inserttotalframenumber\end{beamercolorbox}}}
\setbeamersize{text margin left=12mm,text margin right=12mm}
\begin{document}
\begin{frame}[t]{}
\vspace{14mm}\textcolor{InsperRed}{\bfseries AULA 01}\par\vspace{4mm}{\LARGE\bfseries Fundamentos dos Large Language Models}\par\vspace{6mm}{\large Marketing Analítico}\par\vspace{12mm}{\color{gray}Paulo C. Marques F. · 2027-1}
\end{frame}

\begin{frame}[t]{Objetivos e verificação}
\small
\begin{itemize}
\item Compreender o mecanismo de geração de linguagem e distingui-lo de classificação, cálculo e consulta a bases de dados. Iniciar a prática de especificação e verificação de software por prompts no Codex (interface gráfica).
\item Resolver manualmente os exemplos do capítulo, identificando hipóteses e consequências.
\item Especificar a aplicação no Codex gráfico e confrontar ao menos um resultado com evidências.
\end{itemize}
\end{frame}

\begin{frame}[t]{Organização do encontro}
\small
\begin{itemize}
\item Apresentação dos fundamentos e resolução de problemas quantitativos no papel.
\item Laboratório supervisionado e testes produzidos por prompts na interface gráfica do Codex, sem terminal.
\item Atividade individual com dados diferentes e depósito no GitHub antes do quiz seguinte.
\end{itemize}
\end{frame}

\begin{frame}[t]{Linguagem como sequência de tokens}
\small
\begin{itemize}
\item Um Large Language Model (LLM) autorregressivo recebe uma sequência de tokens e estima uma distribuição de probabilidade para o token seguinte. Token não é sinônimo de palavra: uma palavra pode ser representada por…
\item Consequentemente, o custo de contexto depende da tokenização efetiva.
\item O limite de contexto inclui instruções, histórico, documentos e a própria pergunta, e não apenas o texto visível da última mensagem.
\item Denotemos por a sequência tokenizada.
\end{itemize}
\end{frame}

\begin{frame}[t]{Linguagem como sequência de tokens — continuação}
\small
\begin{itemize}
\item A regra da cadeia da probabilidade fornece Essa fatoração não afirma que o modelo consulta fatos externos; descreve o problema probabilístico que a modelagem resolve.
\item O treinamento ajusta parâmetros a partir de sequências; na inferência, o modelo aplica os parâmetros já aprendidos a um contexto particular.
\item A aplicação de negócios é uma terceira camada: especifica como converter textos em resultados que possam ser verificados.
\end{itemize}
\end{frame}

\begin{frame}[t]{Linguagem como sequência de tokens — formulação matemática}
\small
\vspace{3mm}\begin{center}\small\[P(x_1,\ldots,x_n)=\prod_{i=1}^{n}P(x_i\mid x_1,\ldots,x_{i-1}).\]\end{center}
\vspace{3mm}\small Essa fatoração não afirma que o modelo consulta fatos externos; descreve o problema probabilístico que a modelagem resolve.
\end{frame}

\begin{frame}[t]{Decodificação e variabilidade}
\small
\begin{itemize}
\item Se os escores de três tokens forem , a distribuição normalizada por softmax à temperatura é A escolha greedy seleciona o maior ; a amostragem pode selecionar alternativas menos prováveis.
\item Temperatura controla a dispersão da distribuição, não a veracidade.
\item Temperatura reduzida não equivale a uma verificação de fontes.
\item O exemplo da tabela separa probabilidade de correção factual.
\end{itemize}
\end{frame}

\begin{frame}[t]{Decodificação e variabilidade — continuação}
\small
\begin{itemize}
\item A chance de a amostragem retornar uma alternativa a ``mantido'' é .
\item Isso pode ser calculado sem executar nenhum modelo.
\end{itemize}
\end{frame}

\begin{frame}[t]{Decodificação e variabilidade — formulação matemática}
\small
\vspace{3mm}\begin{center}\small\[p_i=\frac{\exp(z_i/\tau)}{\sum_j\exp(z_j/\tau)}.\]\end{center}
\vspace{3mm}\small A escolha greedy seleciona o maior ; a amostragem pode selecionar alternativas menos prováveis.
\end{frame}

\begin{frame}[t]{Atenção como combinação condicionada ao contexto}
\small
\begin{itemize}
\item No mecanismo de scaled dot-product attention, para uma matriz de representações , definimos , e .
\item Uma cabeça de atenção calcula em que impede que uma posição causal acesse tokens futuros.
\item O produto entre queries e keys produz escores de compatibilidade; os pesos após softmax combinam os valores.
\item A operação não é uma busca na internet, e pesos de atenção elevados não constituem prova de influência causal.
\end{itemize}
\end{frame}

\begin{frame}[t]{Atenção como combinação condicionada ao contexto — continuação}
\small
\begin{itemize}
\item Em um exemplo escalar com escores não normalizados e e valores e , os pesos são e , de modo que a combinação é .
\end{itemize}
\end{frame}

\begin{frame}[t]{Atenção como combinação condicionada ao contexto — formulação matemática}
\small
\vspace{3mm}\begin{center}\small\[\operatorname{Attention}(Q,K,V)=\operatorname{softmax}\!\left(\frac{QK^\top}{\sqrt{d_k}}+M\right)V,\]\end{center}
\vspace{3mm}\small em que impede que uma posição causal acesse tokens futuros.
\end{frame}

\begin{frame}[t]{Evidência, extração e inferência}
\small
\begin{itemize}
\item Considere a avaliação ``A embalagem chegou intacta, mas o produto falhou duas vezes''.
\item São sustentadas as afirmações sobre embalagem e falhas; não são sustentadas afirmações sobre fraude, reembolso ou intenção de recompra.
\item Um painel de Marketing Analítico deve distinguir observação textual, classificação sugerida e inferência do analista.
\item Uma citação literal verificável é mais forte do que uma descrição plausível sem rastreabilidade.
\end{itemize}
\end{frame}

\begin{frame}[t]{Evidência, extração e inferência — continuação}
\small
\begin{itemize}
\item Identificadores de registros permitem voltar à fonte; casos inconclusivos não devem ser convertidos artificialmente em categorias certas. Exemplo resolvido. Uma janela comporta 16 tokens.
\item Instruções consomem 5, documentos recuperados 6 e a pergunta 7.
\item Remover indiscriminadamente os dois tokens mais antigos pode eliminar uma negação ou referência.
\item A solução deve preservar o significado, não apenas reduzir a contagem.
\end{itemize}
\end{frame}

\begin{frame}[t]{Exercício para resolver no papel 1}
\small
\small\textbf{1.1 A próxima palavra do comunicado}\par\vspace{4mm}Um gerador deve completar a expressão `O prazo será''. Para os tokens mantido'', ampliado'' e revisto'', as probabilidades são 0,50, 0,35 e 0,15. (a) Indique a saída gulosa; (b) determine a chance de amostrar algo diferente de `mantido''; (c) explique por que repetir a geração pode produzir resultados distintos sem mudar os dados.
\end{frame}

\begin{frame}[t]{Exercício para resolver no papel 2}
\small
\small\textbf{1.2 Duas etapas, uma notícia}\par\vspace{4mm}Para o início `A empresa'', suponha e . Calcule a probabilidade da sequência de dois tokens. Um colega multiplica 0,4 pela frequência incondicional da palavra `preços'', igual a 0,3. Identifique seu erro.
\end{frame}

\begin{frame}[t]{Laboratório e caso aplicado}
\small
\begin{itemize}
\item Área de aplicação: Marketing Analítico.
\item Conferir os 36 registros e comparar manualmente três comentários: elogio inequívoco, crítica e caso ambíguo.
\item Especificar no Codex (interface gráfica) um painel com contagens por marca, temas propostos, frases de apoio e campo ``evidência insuficiente''.
\item Solicitar testes de existência de IDs, contagens e correspondência literal entre as citações e os textos de origem; inspecionar dois casos problemáticos.
\end{itemize}
\end{frame}

\begin{frame}[t]{Prompt de desenvolvimento — exemplo}
\small
\small\textbf{Instrução fornecida ao Codex:}\par\vspace{5mm}Crie um painel local em Python para ler os 36 comentários fictícios, organizar temas com exemplos literais e registrar afirmações incertas. Gere testes das contagens e das citações. Construa e execute tudo sem API convencional; quando necessário, integre o Codex Python SDK apenas após validar o ambiente.
\end{frame}

\begin{frame}[t]{Critérios de aceitação e validação}
\small
\begin{itemize}
\item Nenhuma citação inexistente; 36 IDs preservados; informação ausente nunca inventada; pelo menos um caso classificado como inconclusivo; registro de duas falhas efetivamente encontradas nos testes.
\item Inspecionar argumentos, saídas e rastros das ferramentas quando houver chamadas externas.
\item Comparar uma saída com a resolução manual do exemplo conceitual.
\end{itemize}
\end{frame}

\begin{frame}[t]{Atividade individual em casa}
\small
\begin{itemize}
\item Novo arquivo com 24 avaliações fictícias de acessórios para bicicleta, duas marcas e três comentários contraditórios; incluir notas e IDs. O arquivo não reaproveita textos do laboratório.
\item Resolver as questões de interpretação, com cálculos e evidência documental.
\item Depositar no GitHub os arquivos solicitados; a entrega habilita o quiz, sem valer nota.
\end{itemize}
\end{frame}

\begin{frame}[t]{Preparação para a avaliação escrita}
\small
\begin{itemize}
\item Exercícios do caderno: 1.1 A próxima palavra do comunicado; 1.2 Duas etapas, uma notícia; 1.3 Atenção numa reclamação.
\item Formalizar o raciocínio e interpretar resultados; não será exigida escrita manual de código.
\item No quiz seguinte serão avaliados conceitos desta aula e da atividade feita em casa.
\end{itemize}
\end{frame}

\end{document}
